\section{Related Work}
\label{sec:related}

Our contribution sits at the confluence of three literatures that have, until recently, evolved in relative isolation: the theory of inference-time reward optimization as exponential tilting, the empirical study of in-context learning (ICL) in audio language models, and the geometric characterization of what frozen speech encoders represent. We review each in turn, emphasizing not a catalogue of methods but the structural tensions that motivate our framework: that training-free RL is a search over a fixed action space \cZ whose efficacy is bounded by the support of the base distribution \qo, and that this boundedness manifests differently for the embedding operator \cZA than for the generative operator \cZB.

\subsection{Inference-Time / Training-Free RL and the Tilting View}
\label{sec:rw-tilting}

The unifying abstraction behind training-free RL is that a reward-guided search distribution \qz is, in its optimal form, an exponential tilt of the base policy: the KL-regularized objective \(\Fobj{q}=\Exp_{q}[\Rwd]-\beta^{-1}\KL{q}{\qo}\) is maximized by the Gibbs measure \(\qs(z)\propto\qo(z)\exp(\beta\Rwd(z))\), with partition function \Zpart \citep{korbak2022kl}. This identity recasts a spectrum of seemingly disparate inference-time procedures as approximations to the same tilt. Best-of-$N$ (BoN) sampling is the most studied instance: \citet{beirami2024bon} establish that the KL divergence between the BoN distribution and \qo grows as \(\log N - (N{-}1)/N\), giving a closed-form exchange rate between compute and the distributional shift purchased. \citet{verdun2025softbon} sharpen this with a soft, temperature-controlled BoN whose suboptimality decays as \(O(1/N)\), while \citet{aminian2025smoothing} recast BoN through a smoothing lens, bounding its regret against the Gibbs optimum and exposing when additional samples cease to help. The crucial—and for our purposes load-bearing—observation across this line is that \emph{every} tilt reweights existing samples: the support of \qz is contained in the support of \qo. No amount of $N$ creates mass where \qo assigns none. Minimum Bayes risk (MBR) decoding \citep{ichihara2025mbr}, including its specialization to ASR \citep{mbrasr2025}, and controlled decoding \citep{mudgal2024controlled} inherit the same ceiling: they are consensus or value-guided reweightings, powerful when the correct hypothesis is already sampled with non-negligible probability and inert otherwise.

A more recent strand attempts to extract reward signal without any labels at all. TTRL \citep{zuo2025ttrl} and TPO \citep{li2025tpo} use self-consistency or majority vote as a surrogate reward at test time, and JitRL \citep{jitrl2026} performs gradient-free credit assignment over a frozen policy. These methods are attractive precisely because they touch no weights, but they sharpen rather than dissolve the support constraint: a self-generated reward can only amplify modes the model already prefers, which is why majority-vote rewards can entrench confident errors. The over-optimization literature quantifies the resulting hazard. \citet{gao2023overopt} show that reward and true objective diverge as KL grows, following a predictable scaling law; \citet{khalaf2025hedgetune} characterize the optimal effective $N^\ast$ (the Best-of-Poisson regime) beyond which tilting harms. The over-optimization risk is governed by the reward spread \spread and the realized gain \gain: when the proxy reward is mis-specified, larger \gain trades calibrated base knowledge for spurious peaks.

These dynamics intersect the ongoing RLVR boundary debate, which directly concerns whether inference-time or verifiable-reward optimization can expand a model's capabilities or merely re-allocate them. \citet{rlvrincentive2025} argue that RLVR predominantly sharpens the base distribution rather than adding skills, a position echoed by the finding that even spurious or random rewards can improve benchmark scores by exploiting format and prior structure \citep{spuriousrewards2025}. The countervailing evidence that RLVR can induce genuinely new skills \citep{newskills2026} does not contradict the tilting view so much as locate its exception: new behavior emerges only when the action space \cZ is rich enough—e.g., multi-step generative composition—that the relevant trajectory had non-zero base mass that single-shot sampling rarely surfaced. This distinction—reweighting versus genuine expansion—is exactly the axis along which we separate \cZA from \cZB, and the tilting formalism makes the separation precise rather than rhetorical.

\subsection{In-Context Learning in Audio LLMs and the Two Model Classes}
\label{sec:rw-icl}

If training-free RL reweights the base distribution, the natural question is whether in-context demonstrations can \emph{move} it. The text-domain ICL literature offers a sobering prior. \citet{min2022rethinking} demonstrate that demonstrations are largely label-insensitive: corrupting ground-truth labels barely degrades performance, implying that demonstrations specify the task and output format rather than teaching the input-output mapping. \citet{alice2026} reach a congruent conclusion for instruction-following demonstrations, which fix format but not accuracy. ICL is moreover scale-gated—its more sophisticated, override-the-prior behaviors emerge only in larger models \citep{wei2023larger}—a property that carries into audio, where few-shot competence appears as an emergent phenomenon in larger audio LLMs \citep{mimoaudio2025}. The function- and task-vector view \citep{todd2024functionvectors,taskvectors2025} explains the mechanism: demonstrations compress to a latent task direction that conditions decoding without altering the underlying perceptual representation. For a generative operator \cZB, this means demonstrations can re-route \emph{how} content is verbalized, but they do not re-open perceptual channels the encoder has discarded.

This perceptual ceiling is starkly visible in audio. A growing body of evidence shows that audio LLMs frequently ``read rather than listen'': \citet{voxparadox2026} document that models answer paralinguistic questions from textual priors and transcript content rather than the acoustic signal, and \citet{listenortranscribe2025} find that performance often depends on whether the model transcribes versus genuinely attends to non-lexical cues. In-context demonstrations do little to repair this when the requisite acoustic information has already been pooled away upstream—consistent with the broader finding that few-shot prompting fixes format, not the missing percept. SER is a partial and instructive exception: \citet{sericl2025} report that few-shot demonstrations \emph{do} improve speech emotion recognition, plausibly because emotion is one of the paralinguistic factors a generative omni model retains with enough base probability for tilting and ICL to surface. The pattern—ICL helping for retained factors, failing for discarded ones—mirrors the support constraint of \Cref{sec:rw-tilting} and foreshadows our operator-dependent treatment.

The contrast with embedding models is decisive for our two-class taxonomy. Text embedders do not acquire ICL for free: \citet{bgeenicl2024} show that exploiting in-context examples in an embedder requires dedicated training, and instruction-conditioned embedders such as INSTRUCTOR \citep{su2023instructor} must be explicitly trained to follow task instructions at encode time. An off-the-shelf embedding operator \cZA therefore exposes \emph{no} prompt-conditioned degree of freedom analogous to \cZB's in-context steering: its action space is the choice of pooling, layer, and query, not a re-conditioning of the forward pass. This asymmetry—\cZB admits demonstration-driven tilts, \cZA does not—is, in our framework, not an implementation detail but a defining structural difference between the two model classes, and it dictates which training-free interventions are even available for each.

\subsection{Capability Map, Embedding Geometry, and Probing}
\label{sec:rw-geometry}

Whether a frozen model \emph{can} be tilted toward a task ultimately depends on whether the relevant information survives in its representations—a question the probing and embedding-geometry literatures answer with unusual clarity. Contrastive and alignment objectives are now understood to actively discard nuisance variation: \citet{xiao2020whatnot} show that contrastive learning suppresses factors not aligned with the training augmentations, and \citet{wang2020alignuniform} formalize the alignment–uniformity trade-off that drives semantically similar inputs together while spreading the hypersphere uniformly. For a content-oriented embedder, speaker identity, emotion, and channel are precisely the nuisances such objectives are designed to remove. This is the geometric root of why an embedding operator \cZA, however cleverly searched, cannot recover a factor its training objective erased: the base distribution \qo over \cZA simply has no mass on the discarded direction.

Pooling compounds the loss. Mean- or last-token pooling imposes an information bottleneck that NV-Embed's learned latent-attention pooling was introduced to mitigate \citep{nvembed2024}, and the speaker-verification community has long known that attentive statistics pooling \citep{attnstatspool2018} and ECAPA-TDNN's multi-scale aggregation \citep{ecapatdnn2020} are necessary precisely because naive temporal averaging destroys speaker-discriminative structure. The implication is that the \emph{same} encoder can be near-blind or highly informative for a factor depending solely on the read-out, which is exactly the degree of freedom a training-free search over \cZA should exploit. Layer-wise probing makes this concrete: \citet{pasad2021layerwise} chart an autoencoder-like trajectory across wav2vec~2.0 layers in which phonetic and word-level information peaks at intermediate depths, and \citet{spkattrprobing2025} show speaker attributes are distributed non-uniformly across layers. Capability is thus not a scalar property of a model but a function of (layer, pooling, query)—the very coordinates of our action space \cZA.

Crucially, evidence that paralinguistic content \emph{persists} even when unused is now substantial. \citet{frozenllmparaling2024} demonstrate that a frozen LLM, given speech features, perceives paralinguistic attributes it never verbalizes; \citet{emotionneurons2026} localize emotion-selective units inside such models. These findings establish the optimistic precondition for our flagship study: the information needed to disentangle content, speaker, emotion, and intent is present in the frozen omni model's hidden states, even when the default forward pass—or a generative decode \cZB—suppresses it. The task is therefore one of \emph{activation} via search over read-outs, not of teaching. The SUPERB and Dynamic-SUPERB taxonomies \citep{superb2021,dynsuperb2023} supply the benchmark scaffolding for evaluating such activation across content, speaker, semantic, and paralinguistic axes, and we adopt their factorization to structure our capability map. What this literature lacks—and what our tilting formalism supplies—is a principled account of \emph{which} of these latent capabilities are reachable by reweighting a fixed \qo and which demand a different action space altogether.